\chapter{對稱性、簡併與角動量階梯}
\providecommand{\ket}[1]{\left\lvert#1\right\rangle}

\section{對稱性與簡併能階}

對稱性是一種不改變系統物理描述的操作。它不會改變哈密頓量、Schr\"odinger方程或可能的能量。若以么正算符\(S\)表示此操作，不變性便意味著
\begin{equation}
S H S^\dagger=H,
\end{equation}
或等價地
\begin{equation}
[S,H]=0.
\end{equation}

平移就是一個例子。將原子從一處移到另一處，原子的能量不會改變。這對所有原子態都成立，不只對基態如此。旋轉是另一個例子，也是本節的主要主題。

交換兩個全同粒子也是一種對稱性。把兩個電子稱為電子一和電子二，或稱為 Harry 和 Fred，都不會使它們變得不同。交換標記不會改變原子的性質。

晶體還有一種額外的平移對稱性。將整個晶體在空間中平移，是一般的平移。然而，對無限晶體而言，平移一個晶格間距後，物理排列也會保持不變。等價地說，將每個格點的標號都平移一格，就能重新標記格點。有限晶體的邊界會使這個說法不再嚴格成立；但在很長的晶體中，遠離兩端之處仍可視為具有這項近似對稱性。

空間的均勻性---也就是每個位置都和其他位置一樣---就是平移對稱性。空間的各向同性---也就是每個方向都和其他方向一樣---就是旋轉對稱性。

簡併與對稱性有關，但兩者並不相同。若不只一個態具有完全相同的能量，這個能階便是簡併的。對於具有離散譜的一般哈密頓量，兩個互不相關的能階恰好完全相等，會是驚人的巧合。即使兩者相差不到\(10^{15}\)，仍不代表完全相等。精確簡併通常暗示著對稱性。

對稱性有時會導致簡併，但並非總是如此。第一個例子會說明對稱性為何可能導致簡併，也會說明為何單靠一種對稱性有時還不夠。

\section{圓上的粒子}

考慮一顆被限制在圓形軌道上運動的珠子。它的位置可以用笛卡兒座標\(x\)和\(y\)描述，但只用一個角座標就足夠：
\begin{equation}
\psi=\psi(\theta).
\end{equation}
圓上的位置機率密度為\(\psi^*(\theta)\psi(\theta)\)。

逆時針主動旋轉一個小角度\(\epsilon\)後，波函數變為
\begin{equation}
\psi(\theta)\longrightarrow\psi(\theta-\epsilon).
\end{equation}
因此
\begin{align}
\delta\psi(\theta)
&=\psi(\theta-\epsilon)-\psi(\theta)\\
&=-\epsilon\frac{d\psi}{d\theta}
+\mathcal O(\epsilon^2).
\end{align}

算符
\begin{equation}
L_z=-i\hbar\frac{d}{d\theta}
\end{equation}
作用在週期波函數上時是厄米算符。它是
\begin{equation}
p_x=-i\hbar\frac{d}{dx}.
\end{equation}
在角座標中的類比。因子\(i\)對厄米性至關重要。以\(L_z\)表示，
\begin{equation}
\delta\psi
=-\frac{i\epsilon}{\hbar}L_z\psi
+\mathcal O(\epsilon^2),
\end{equation}
所以無窮小旋轉算符為
\begin{equation}
R(\epsilon)
=1-\frac{i\epsilon}{\hbar}L_z
+\mathcal O(\epsilon^2).
\end{equation}
這表示\(L_z\)是繞垂直於圓面的軸旋轉時的生成元。

\begin{classroomqa}
\lecturetimestamp{00:11:19}
\audiencequestion{涉及\(i\)的其中一個因子，符號是否應該相反？}
\lecturerresponse{檢查符號後，就得到上面的關係。按照約定\(\psi(\theta)\to\psi(\theta-\epsilon)\)，變化量是\(-\epsilon\,d\psi/d\theta\)，而這等於\(-i\epsilon L_z\psi/\hbar\)。}
\end{classroomqa}

現在求角動量的本徵態：
\begin{equation}
L_z\psi=M\psi.
\end{equation}
代入微分算符後得到
\begin{equation}
-i\hbar\frac{d\psi}{d\theta}=M\psi,
\end{equation}
其解的形式為
\begin{equation}
\psi(\theta)=C\exp\left(\frac{iM\theta}{\hbar}\right).
\end{equation}

這裡還有一個在無限直線上的運動不會遇到的額外條件。\(\theta=0\)和\(\theta=2\pi\)是同一個點，因此單值的軌道波函數必須滿足
\begin{equation}
\psi(\theta+2\pi)=\psi(\theta).
\end{equation}
對於上面的本徵函數，這表示
\begin{equation}
\exp\left(\frac{2\pi iM}{\hbar}\right)=1.
\end{equation}
因此，
\begin{equation}
M=m\hbar,
\qquad
m\in\mathbb Z.
\end{equation}
圓上的軌道角動量可能值為
\begin{equation}
\ldots,-2\hbar,-\hbar,0,\hbar,2\hbar,\ldots.
\end{equation}

角動量為\(m\hbar\)的態，是否必須與角動量為\(-m\hbar\)的態具有相同能量？正、負\(m\)代表環流方向相反。問題是，單靠旋轉對稱性，是否就能要求這兩種運動具有相同能量？

\begin{classroomqa}
\lecturetimestamp{00:22:51}
\audiencequestion{採用\(\psi(\theta-\epsilon)\)這個形式，是否表示粒子正順時針運動？}
\lecturerresponse{不是。這個約定只決定哪個方向稱為正向旋轉。它定義角位置和角動量的正負號，卻不決定粒子實際運動的方向。}
\end{classroomqa}

垂直於圓面方向的磁場提供了一個反例。它可以使正、負角動量對應的能量朝相反方向移動，因而使兩者的能量分裂。能量位移的正負取決於電荷與方向的約定；此處重要的是能量分裂這件事。

這種磁場仍保有繞垂直軸的旋轉對稱性。因此，它顯示單靠旋轉對稱性並不能要求
\begin{equation}
E(m)=E(-m).
\end{equation}
若沒有能量分裂，\(m=1\)和\(m=-1\)可以形成一對簡併態，\(m=2\)和\(m=-2\)也可以。態\(m=0\)則沒有這種不同的配對態。

\section{\texorpdfstring{\(m,-m\)}{m, -m}對與反射對稱性}

再加入一種離散反射對稱性。令\(\mathcal M\)表示對圓的一條直徑作鏡射：
\begin{equation}
(\mathcal M\psi)(\theta)=\psi(-\theta).
\end{equation}
它會反轉環流方向：
\begin{equation}
\mathcal M\psi_m=\psi_{-m}.
\end{equation}

若反射是對稱操作，則
\begin{equation}
[\mathcal M,H]=0.
\end{equation}
假設
\begin{equation}
H\ket{m}=E_m\ket{m}.
\end{equation}
作用\(\mathcal M\)後得到
\begin{align}
H\mathcal M\ket{m}
&=\mathcal M H\ket{m}\\
&=E_m\mathcal M\ket{m}.
\end{align}
由於\(\mathcal M\ket{m}\)是角動量為\(-m\hbar\)的態，
\begin{equation}
E_{-m}=E_m.
\end{equation}
旋轉對稱性使角動量守恆；反射對稱性則把兩個相反的角動量值聯繫起來。兩者合在一起，便要求這一對態簡併。

\begin{classroomqa}
\lecturetimestamp{00:29:54}
\audiencequestion{既然\(m\)和\(-m\)態具有相同的機率密度，這是否已足以顯示它們的能量相等？}
\lecturerresponse{不。波函數\(e^{im\theta}\)和\(e^{-im\theta}\)本身不同，儘管它們的絕對值平方相同。角動量本徵值方程涉及波函數本身，而不只是\(\psi^*\psi\)。量子力學需要的是振幅，不能只看機率。}
\end{classroomqa}

磁場保有繞其軸的旋轉對稱性，卻破壞了前面所用的反射對稱性。

\begin{classroomqa}
\lecturetimestamp{00:33:05}
\audiencequestion{反射對稱性會使磁場反向嗎？}
\lecturerresponse{會。磁場是軸向量。將指向紙面內的磁場作鏡射，會得到指向紙面外的磁場。等價地說，電流環及其鏡像中的電流方向相反，因此產生方向相反的磁場。}
\end{classroomqa}

旋轉與反射不對易，這點可以直接看出。作用在角動量本徵函數上時，
\begin{align}
\mathcal M L_z\psi_m
&=m\hbar\,\psi_{-m},\\
L_z\mathcal M\psi_m
&=-m\hbar\,\psi_{-m}.
\end{align}
因此
\begin{equation}
[\mathcal M,L_z]\psi_m
=2m\hbar\,\psi_{-m},
\end{equation}
當\(m\neq0\)時，結果不為零。先旋轉再反射，與先反射再旋轉，並非同一個操作。

當兩種變換都是對稱操作時，兩個算符都與\(H\)對易，但兩者彼此不對易。這就是導致簡併的模式。

\begin{classroomqa}
\lecturetimestamp{00:40:28}
\audiencequestion{任意一對不對易的海森堡算符，例如\(x\)和\(p\)，都會產生簡併嗎？}
\lecturerresponse{一般而言不會。這些算符還必須是對稱性，也就是必須與哈密頓量對易。位置和動量通常不會同時守恆。若兩者確實都與\(H\)對易，那麼它們的不對易性就會意味著簡併。}
\end{classroomqa}

\section{對稱性的對易子代數}

假設\(A\)和\(B\)是兩個連續對稱性的厄米生成元：
\begin{equation}
[A,H]=0,
\qquad
[B,H]=0.
\end{equation}
這些方程並未要求\(A\)和\(B\)彼此對易。

對易子\([A,B]\)是反厄米的，因此定義一個厄米算符
\begin{equation}
C=i[A,B].
\end{equation}
此處數值因子和整體正負號並不重要。關鍵是\(C\)也與哈密頓量對易。確實，
\begin{equation}
[AB,H]
=A[B,H]+[A,H]B=0,
\end{equation}
同樣地
\begin{equation}
[BA,H]=0.
\end{equation}
因此
\begin{equation}
[C,H]
=i\bigl([AB,H]-[BA,H]\bigr)=0.
\end{equation}

新的生成元\(C\)可能獨立於\(A\)和\(B\)，也可能只是已知生成元的線性組合。若它是新的，就將它加入列表，再計算更多對易子。這個過程可能無限延續，但通常會在有限個生成元之後閉合：所有後續對易子都是現有生成元的線性組合。這個閉合的系統稱為對易子代數或李代數。

有限的對稱變換構成一個群；它們的無窮小生成元則構成相應的李代數。

簡併之所以重要，是因為它能整理能譜。若兩個態能量相同，一個態就不能藉由發射能量等於兩態能量差的光子而衰變至另一個態。在原子光譜學中，能階的排列方式及其簡併情形決定可能出現的譜線。同樣的道理也適用於激發態強子及其躍遷。

\begin{classroomqa}
\lecturetimestamp{00:49:07}
\audiencequestion{這些對易子本身是對稱群的元素嗎？}
\lecturerresponse{這裡討論的是生成元。生成元及其對易子構成對易子代數，也就是李代數。由生成元構成的有限對稱操作才是群元素。}
\end{classroomqa}

若所有對稱性生成元彼此對易，這種對稱性稱為阿貝爾對稱性。若有些生成元的對易子不為零，則稱為非阿貝爾對稱性。

\begin{classroomqa}
\lecturetimestamp{00:50:25}
\audiencequestion{生成元的根本定義是什麼？}
\lecturerresponse{取一個非常接近恆等變換的么正對稱變換，將它寫成\(1-i\epsilon G/\hbar+\cdots\)。厄米算符\(G\)就是這個變換的生成元。重複相應的無窮小變換，就能構成有限變換。}
\end{classroomqa}

\begin{classroomqa}
\lecturetimestamp{00:52:08}
\audiencequestion{對稱性生成元的代數是什麼意思？}
\lecturerresponse{它是由一組生成元構成的集合；這些生成元可以相加、相減、乘以一般數值，也可以彼此取對易子。在這個代數中，對易子扮演乘法的角色；封閉性表示這些運算不會產生集合以外的生成元。}
\end{classroomqa}

\section{繞不同軸旋轉}

旋轉是非阿貝爾對稱性最重要的例子。先繞\(x\)軸旋轉，再繞\(y\)軸旋轉，一般不等同於以相反順序進行這兩次旋轉。兩種順序之間的差異本身與繞第三軸旋轉有關。

\begin{classroomqa}
\lecturetimestamp{00:54:22}
\audiencequestion{有限旋轉不對易，但無窮小旋轉是否對易？}
\lecturerresponse{兩者的差異出現在兩個小旋轉角的二階項。這就足夠了：提出這兩個小角度因子後，生成元本身的對易子仍不為零。}
\end{classroomqa}

在計算量子對易子之前，先看看對應的經典圖像。考慮一個圓形行星軌道，其角動量垂直於軌道平面。將整個軌道繞\(y\)軸旋轉，會得到另一個能量完全相同的允許軌道。不過，這個軌道的角動量向量現在具有不同的\(x\)分量。

因此，旋轉不變性會產生能量相同、角動量分量不同的組態。在量子力學中，這表現為繞不同軸旋轉的生成元彼此不對易。

\section{笛卡兒座標中的角動量}

先考慮一個在\(xy\)平面內運動、但不受限於圓周的粒子。作微小的逆時針旋轉時，
\begin{equation}
\delta x=-\epsilon y,
\qquad
\delta y=\epsilon x.
\end{equation}
對於主動旋轉後的態，變換後的波函數要在逆旋轉後的座標處取值：
\begin{equation}
\psi_\epsilon(x,y)
=\psi(x+\epsilon y,y-\epsilon x).
\end{equation}
因此，
\begin{align}
\delta\psi
&=\epsilon\left(
y\frac{\partial}{\partial x}
-x\frac{\partial}{\partial y}
\right)\psi
+\mathcal O(\epsilon^2)\\
&=-\frac{i\epsilon}{\hbar}
\left(xp_y-yp_x\right)\psi
+\mathcal O(\epsilon^2).
\end{align}
所以，繞\(z\)軸旋轉的生成元為
\begin{equation}
L_z=xp_y-yp_x.
\end{equation}
這就是熟知的經典向量
\begin{equation}
\mathbf L=\mathbf r\times\mathbf p.
\end{equation}
的\(z\)分量。它不只是力學角動量；在量子力學中，它也是使波函數旋轉的算符。\footnote{採用主動旋轉約定，是為了與先前的變換\(\psi(\theta)\mapsto\psi(\theta-\epsilon)\)以及口頭推導中經過符號確認後得到的結果\(L_z=xp_y-yp_x\)一致。}

對\(x\to y\to z\to x\)作循環置換，就能得到其他分量。\footnote{\(L_x\)和\(L_y\)的公式，是依照講課中明確指定的循環置換\(x\to y\to z\to x\)，由\(L_z=xp_y-yp_x\)推得的。}
\begin{align}
L_x&=yp_z-zp_y,\\
L_y&=zp_x-xp_z,\\
L_z&=xp_y-yp_x.
\end{align}

對於旋轉不變的系統，這些算符都與哈密頓量對易：
\begin{equation}
[H,L_x]=[H,L_y]=[H,L_z]=0.
\end{equation}
在中心力場中運動的粒子是標準例子。這些方程就是角動量守恆的量子形式。

\section{角動量李代數}

正則對易關係為
\begin{align}
[x_i,x_j]&=0,\\
[p_i,p_j]&=0,\\
[x_i,p_j]&=i\hbar\,\delta_{ij}.
\end{align}
座標與不同方向的動量分量彼此對易。因此，例如\(x\)與\(p_y\)對易，而\(x\)與\(p_x\)不對易。

現在計算
\begin{equation}
[L_x,L_y]
=
[yp_z-zp_y,\;zp_x-xp_z].
\end{equation}
大多數項都會消失。僅有兩個非零項：
\begin{align}
[yp_z,zp_x]&=-i\hbar\,yp_x,\\
[zp_y,xp_z]&=i\hbar\,xp_y.
\end{align}
因此
\begin{align}
[L_x,L_y]
&=i\hbar\left(xp_y-yp_x\right)\\
&=i\hbar L_z.
\end{align}
對座標標記作循環置換，可得
\begin{align}
[L_x,L_y]&=i\hbar L_z,\\
[L_y,L_z]&=i\hbar L_x,\\
[L_z,L_x]&=i\hbar L_y.
\end{align}

\begin{classroomqa}
\lecturetimestamp{01:14:50}
\audiencequestion{當對易子產生\(L_z\)時，我們是否仍須證明\(L_z\)不是\(L_x\)和\(L_y\)的線性組合？}
\lecturerresponse{是的。必須證明它們彼此獨立，不能直接假定如此。此處不進一步作直接的代數證明。}
\end{classroomqa}

上面列出的對易關係在\(L_x\)、\(L_y\)和\(L_z\)所張成的空間中閉合：將三者都納入後，後續的對易運算不會產生此空間以外的算符。這就是角動量李代數。再加上
\begin{equation}
[H,L_i]=0,
\qquad i=x,y,z,
\end{equation}
這些關係便包含了能譜的大量資訊，無須直接求解角向Schr\"odinger方程。

量子諧振子也採用類似的代數方法，以產生算符和湮滅算符生成其能譜。在這裡，類似的算符改變的是角動量的\(z\)分量，而不是能量。

選擇\(z\)軸只是選擇一種表象。物理上並沒有理由認為\(z\)比\(x\)或\(y\)更特殊，但選定一個分量可得到方便的基底。令
\begin{equation}
L_z\ket{m}=m\hbar\ket{m}.
\end{equation}
我們要找出\(m\)可能取哪些值，以及不同\(m\)的態之間有何關係。

\begin{classroomqa}
\lecturetimestamp{01:24:02}
\audiencequestion{你將這個構造和諧振子比較時，指的是量子諧振子嗎？上一堂課並沒有講到它。}
\lecturerresponse{是的，比較的對象是量子諧振子。這裡的角動量論證本身完整，不需要先了解諧振子。之後也可以反過來看這個類比：諧振子同樣運用了類似的代數方法。}
\end{classroomqa}

\section{角動量的升降算符}

定義
\begin{equation}
L_+=L_x+iL_y,
\qquad
L_-=L_x-iL_y.
\end{equation}
角動量的對易關係給出
\begin{equation}
[L_z,L_+]=\hbar L_+,
\qquad
[L_z,L_-]=-\hbar L_-.
\end{equation}

將第一個關係式作用在\(L_z\)的本徵態上：
\begin{align}
L_zL_+\ket{m}
&=\left(L_+L_z+[L_z,L_+]\right)\ket{m}\\
&=(m+1)\hbar\,L_+\ket{m}.
\end{align}
除非\(L_+\ket{m}=0\)，否則所得結果是\(L_z\)的另一個本徵態，其本徵值為\((m+1)\hbar\)。同樣地，
\begin{equation}
L_zL_-\ket{m}
=(m-1)\hbar\,L_-\ket{m}.
\end{equation}
因此，\(L_+\)會使\(m\)增加一，而\(L_-\)會使它減少一。

從任一個本徵值出發，代數會產生一個彼此相差整數的序列。這個序列可能無限延伸，也可能因為在上端有
\begin{equation}
L_+\ket{m_{\mathrm{top}}}=0
\end{equation}
以及在下端有
\begin{equation}
L_-\ket{m_{\mathrm{bottom}}}=0
\end{equation}
而終止。

\begin{classroomqa}
\lecturetimestamp{01:35:11}
\audiencequestion{圓的例子不是已經要求角動量必須是整數了嗎？}
\lecturerresponse{那個論證證明，由單值的位置波函數表示的軌道角動量必須取整數值。抽象的角動量代數給出的結果較寬鬆：階梯上可以是整數，也可以是半整數。}
\end{classroomqa}

如果階梯終止，旋轉不變性會把它的兩端聯繫起來。繞垂直於\(z\)的軸旋轉\(180^\circ\)，會使
\begin{equation}
L_z\longrightarrow-L_z.
\end{equation}
因此，譜必須以零為中心對稱：
\begin{equation}
m_{\mathrm{bottom}}=-m_{\mathrm{top}}.
\end{equation}
由於相鄰值相差一，可能的情形有兩種。階梯上的值可以是整數，
\begin{equation}
\ldots,-2,-1,0,1,2,\ldots,
\end{equation}
也可以是半整數，
\begin{equation}
\ldots,-\frac{5}{2},-\frac{3}{2},-\frac{1}{2},
\frac{1}{2},\frac{3}{2},\frac{5}{2},\ldots,
\end{equation}
其端點則位於升算符或降算符作用後結果為零之處。

圓上的單值波函數論證，為軌道角動量選出整數這種可能性，
\begin{equation}
\mathbf L=\mathbf r\times\mathbf p.
\end{equation}
半整數的可能性也符合物理實況，屬於自旋角動量，稍後會另外討論。

最後，若哈密頓量具有旋轉不變性，階梯上的態就具有相同能量。假設
\begin{equation}
H\ket{m}=E\ket{m}.
\end{equation}
由於每個分量\(L_i\)都與\(H\)對易，\(L_+\)和\(L_-\)也會與它對易。因此
\begin{align}
H L_+\ket{m}
&=L_+H\ket{m}\\
&=E L_+\ket{m}.
\end{align}
只要\(L_+\ket{m}\neq0\)，升階後的態就與原態具有相同能量。同樣的計算也會得到
\begin{equation}
H L_-\ket{m}=E L_-\ket{m}.
\end{equation}
因此，升降算符會生成一族\(L_z\)本徵值不同、能量卻完全相同的態。這就是本講一路探討的非對易對稱性與簡併之間的關聯。

\begin{classroomqa}
\lecturetimestamp{01:44:13}
\audiencequestion{這些階梯態是否對應於選擇不同的旋轉軸？}
\lecturerresponse{它們彼此可透過旋轉聯繫。\(L_x\)或\(L_y\)的本徵向量，都是所選\(L_z\)本徵向量的線性組合。旋轉一個態會改變它的角動量分量，但不改變它的能量。}
\end{classroomqa}

選擇\(L_+\)和\(L_-\)很方便，因為這兩個組合能直接呈現階梯結構。這種技巧在量子力學中一再出現：若對易子代數閉合，升算符和降算符往往能整理本徵值譜。

\begin{classroomqa}
\lecturetimestamp{01:46:33}
\audiencequestion{當\(L_+\)作用在最高態上時會怎樣？最高態如何定義？}
\lecturerresponse{作用結果為零，這就是終止階梯頂端的定義。同樣地，\(L_-\)在底端作用時也會得到零。抽象的階梯不一定會終止；不過，對於稍後要討論的原子角動量態，階梯通常會終止。}
\end{classroomqa}

下一步是進一步發展這套角動量結構，並將它應用於原子，包括氫原子。這裡找到的旋轉多重態解釋了部分簡併模式，還有其他結構需要進一步探討。